\addtocontents{toc}{\protect\Needspace{3\baselineskip}}
\originalchapter{分布函数、稳定性与星系演化}\label{ch:distribution}
\originalsection{从轨道到统计描述}
令 $f(\boldsymbol x,\boldsymbol v,t)\dd^3x\dd^3v$ 为相空间体元中的质量. Hamilton 流保持相空间体积, 没有碰撞时沿轨道 $\dd f/\dd t=0$, 即
\begin{equation}\label{eq:cbe}
\frac{\partial f}{\partial t}+\boldsymbol v\cdot\nabla_xf
-\nabla\Phi\cdot\nabla_vf=0,\qquad
\nabla^2\Phi=4\pi G\int f\dd^3v.
\end{equation}
势与分布必须自洽, 不能给任意 $f$ 再无条件套用另一个势. 对速度积分时要求 $f$ 在大速度处衰减, 边界项消失. 零阶矩给连续方程; 乘 $v_i$ 再积分得
\begin{equation}
\partial_t(\rho\bar v_i)+\partial_j(\rho\overline{v_iv_j})
=-\rho\partial_i\Phi.
\end{equation}
这是 Jeans 方程, 不是一个独立封闭状态方程. 稳态、球对称、无整体径向流时,
$\dd(\rho\sigma_r^2)/\dd r+2\beta\rho\sigma_r^2/r=-\rho GM(<r)/r^2$,
$\beta=1-(\sigma_\theta^2+\sigma_\phi^2)/(2\sigma_r^2)$. 不同 $\beta$ 可产生相近投影视向速度弥散, 形成质量与各向异性的退化.

\originalsection{盘的局部稳定性}
对薄气体盘取轴对称、波长远小于盘尺度的扰动, 局部径向惯性项给 $\kappa^2$, 压强给 $c_s^2k^2$. 薄片的扰动势满足 $\delta\Phi=-2\pi G\delta\Sigma/|k|$, 因而
\begin{equation}
\omega^2=\kappa^2+c_s^2k^2-2\pi G\Sigma|k|.
\end{equation}
对于 $k>0$, 二次式在 $k=\pi G\Sigma/c_s^2$ 最小. 所有局部轴对称模式稳定需 $Q_g=\kappa c_s/(\pi G\Sigma)\ge1$. 恒星盘的无碰撞响应不同, 不能把恒星速度弥散直接代入气体系数并声称同一个精确判据. 这也说明自引力、速度支持和旋转共同决定稳定性, 单凭盘半径无法在未指定模型时判断全局扰动.

\originalsection{恒星族与化学演化}
恒星族综合把初始质量函数 $\xi(m)$ 与单星年龄、组成依赖的光谱 $l_\nu(m,t,Z)$ 卷积, 再对过去的恒星形成率积分. 最亮星未必贡献多数质量, 尘埃还能改变光谱; 因而以光谱恢复历史是有退化的逆问题.

封闭箱模型假设没有流入流出、瞬时回收和充分混合. 令总质量 $M_T$ 固定, 气体质量 $M_g$, 锁入长寿命恒星的增量 $\dd M_s=-\dd M_g$. 金属气体质量的变化为 $\dd(ZM_g)=(y-Z)\dd M_s$, 其中 $y$ 是每锁入单位质量的新生成金属产额. 展开左侧并消去 $\dd M_g$ 得 $M_g\dd Z=y\dd M_s$, 初始 $Z=0$ 时积分为 $Z=-y\ln\mu$, $\mu=M_g/M_T$. 这条对数关系是模型推导, 不是任意星系的经验定律. G 型矮星问题、吸积和流出提示它的局限.

活动星系核的能量由黑洞吸积提供, 大致 $L=\eta\dot Mc^2$. 尘埃环的方向选择可使同一中心发动机在光谱上表现不同, 但所有 AGN 差异并不都由方向解释. 喷流需要磁场、角动量和相对论流动, 本卷使用能量与 Eddington 尺度, 不声称完成磁流体喷流理论.

\originalsection{局部盘色散关系的来源}
在半径 $R_0$ 的小区域, 取扰动径向波数 $k$ 与频率 $\omega$, 背景面密度近似常数. 线性连续式是 $-\ii\omega\delta\Sigma+\ii k\Sigma u_R=0$. 径向与切向 Euler 式为 $-\ii\omega u_R-2\Omega u_\phi=-\ii k(c_s^2\delta\Sigma/\Sigma+\delta\Phi)$, $-\ii\omega u_\phi+\kappa^2u_R/(2\Omega)=0$. 第二式所含的 $\kappa$ 来自差分旋转, 不是任意添上的振荡频率.

无限薄片的扰动 Poisson 方程为 $(\partial_z^2-k^2)\delta\Phi=4\pi G\delta\Sigma\delta(z)$. 两侧衰减解 $\delta\Phi=C\ee^{-|k||z|}$ 的导数跨片跳变 $-2|k|C$, 故 $C=-2\pi G\delta\Sigma/|k|$. 再由连续式消去 $\delta\Sigma$、切向式消去 $u_\phi$, 得正文色散关系. 若盘有限厚度, 引力响应因垂直分布被减弱; 若扰动非轴对称或波长接近盘尺度, 这个局部推导不能判断所有全局棒模或旋臂模.

Jeans 方程也可在一个自足模型中检验. 球对称、各向同性系统 $\rho=C/r^2$, 外部势给恒定圆速 $v_c$, 取常数 $\sigma_r^2$. Jeans 方程左侧为 $-2C\sigma_r^2/r^3$, 右侧为 $-Cv_c^2/r^3$, 因而 $\sigma_r^2=v_c^2/2$. 这是等温球尺度关系; 截断、各向异性或可见恒星作为示踪体时, 同一个比例不必保留. 通过检验例能看清密度梯度如何进入动力学质量估计.
